\nextpage
\section{Что менялось во входном тексте}
\subhead{Учебный корпус}
Из псевдокода курса подготовлены 72 фрагмента. Первый содержит функцию модуля, далее встречаются преобразования грамматик, обходы и проверки условий. Для записи в текстовый файл знак присваивания приведён к \code{:=}, знак «меньше или равно» к \code{<=}, а греческие буквы в именах заменены словами, например \code{alpha}. Список замен приложен к исходникам. Замерялся один и тот же файл без преобразования переводов строк. Его объём 17\,430 байт, в том числе 530 двухбайтовых переводов строки. На пропуски приходится 46,6\,\% байтов. Среди 3\,792 выданных лексем есть все 31 вида, в том числе 1\,750 идентификаторов и 137 чисел.

\subhead{Общие начала слов}
В трёх файлах по 1\,048\,576 байт повторяются \code{abx}, чередующиеся \code{abx} и \code{ifx} или только \code{ifx}. Далее сравниваются слова длины шесть \code{abcde2} и \code{while2}. За каждым словом стоит один пробел. Внутри каждой группы совпадают объём, число, виды и границы лексем. Меняется доля слов, у которых автомат сначала проходит по началу ключевого слова. Несмотря на это, результатом всякий раз служит один идентификатор.

\subhead{Длина идентификатора и числа}
В файл объёмом $2^{20}$ байт помещаются блоки из 2, 8, 32 или 128 букв \code{a} и одного пробела. Остаток в конце заполняется пробелами. Процедуру повторяют с цифрой \code{1}. Меняются длина лексемы и число выдач за проход. Доля пробелов тоже меняется. Пара «буквы~--- цифры» одной длины даёт одинаковые размеры и позиции лексем при разных видах.

\subhead{Пробелы и комментарии}
Первый файл содержит строки из \code{x}, семи пробелов и перевода строки. Во втором через строку семь пробелов заменены на \code{//aaaaa}, в третьем заменены все. Границы всех \code{x} совпадают. В четвёртом файле на каждые 64 имени приходится один длинный комментарий. Как и в третьем файле, на комментарии приходится 448 байт на 64 имени. Теперь позиции отдельных имён изменились. Во всех четырёх случаях $N=1\,179\,648$, $K=131\,072$.

\subhead{Объём при постоянном составе}
Строка \code{if x:=15 then return x end if} вместе с переводом строки занимает 30 байт и даёт 9 лексем. Её повторяют 2\,048, 32\,768 и 524\,288 раз. Каждый шаг увеличивает $N$ и $K$ в 16 раз, а число выдач на один входной байт остаётся равным $9/30$.
